\documentclass[10pt,norsk,a4paper,usenames,dvipsnames]{article}
\usepackage[utf8]{inputenc}
\usepackage[T1]{fontenc}
\usepackage[norsk]{babel}
% - PDF-relatert
\usepackage{hyperref,pdfpages,hypcap,download}
\hypersetup{colorlinks=true,allcolors=.}
\newcommand\fhref[2]{%
	\href{#1}{#2}\footnote{\url{#1}}%
}
% Andre pakker
\usepackage[cm]{fullpage}
\usepackage{parskip,multicol,textcomp,amssymb,graphicx,color,enumitem,cleveref}
% - Korreksjon av fotnoter i seksjoner/overskrifter
\usepackage[stable]{footmisc}
% - Skrifttype
\usepackage[bitstream-charter]{mathdesign}
% - Kommentarer
\usepackage{comment}
\usepackage{epstopdf}
\usepackage[dvipsnames]{xcolor}
\usepackage{xfrac}
% Redefines the first level of enumerated lists.
\renewcommand{\theenumi}{\alph{enumi}}
\renewcommand{\labelenumi}{\theenumi)}

\epstopdfsetup{outdir=./}

\download[vedtekter.pdf]{protokoll/cyb-gf-protokoll-2021-vår.pdf}


\usepackage{setspace}



\title{\huge Ordinær generalforsamling}
\author{\LARGE Cybernetisk Selskab}
\date{19. mai 2022}

\begin{document}


\maketitle



\begin{center}


\includegraphics[width=0.6\textwidth,height=0.6\textheight,keepaspectratio=true]{cyblogoa3.pdf}

\end{center}


\newpage


\tableofcontents

\section{Valg av møteleder}
Sara mediå Vedtatt ved akklamasjon

\section{Valg av referent}
HIlmar Elverhøy Vedtatt ved akklamasjon

\section{Valg av protokollunderskrivere}
Kenneth Andre Sydnes og
Martin Heggem\\
Vedtatt ved akklamasjon

\section{Valg av tellekorps}
Johan Bjurgård og
Ahmed Abdulahi Ahmed\\
Vedtatt ved akklamasjon

\section{Godkjenning av innkalling}
Vedtatt ved akklamasjon

\section{Godkjenning av dagsorden}
Vedtatt ved akklamasjon
\section{Semesterberetninger}
    \subsection{Semesterberetning ved leder}
    Leder beretter.
    \subsection{Semesterberetning ved kjellermogul}
    Kjellermogul beretter.
    
\section{Kasserer orienterer om økonomi}
    Kasserer beretter.\\
    Nesten ferdig regnskap presenteres. Over 1000 medlemskap er rekord. Samlet salg rundt 700 000kr i 2021. kjøpte inn varer til 600 000 kr. Vi har brukt veldig mye gass. brukt 50 000 kr på sosialt.\\
    Det er ikke satt opp nytt budsjett, da dette bør være ses på av styret\\
    Ingen spørsmål/kommentarer

\section{Kontingentfastsettelse}
    Hovedstyret foreslår å holde medlemskontingenten på kr.~50,-.
    
    Gyda forklarer om saken.
    Ingen innlegg til debatt.\\
    Vedtatt ved akklamasjon

\section{Valg}

    
    \subsection{Hovedstyret}
        Man velges inn i hovedstyret for to semestre av gangen.
        Gyda forteller om hovedstyret, og vervene som er til valg.
        
        \subsubsection{Leder}
        Gyda forklarer om hva det vil si å være leder i en forening. Det er et stort verv, og styre retningen til foreningen.
        
        Hanne Hugem - Er utlånsansvarlig og leder i SPF. Vil representere cyb og gjøre en innsats for foreningen\\
        Peter Petersen tidligere teknisk og nåværende nestleder, har innsikt og vil hjelpe foreningen videre
        
        14:  har dere tid
        ja ja
        
        46: har dere evne
        ja ja
        
        44: liker dere pingviner
        Hanne har 2 pingviner hjemme
        Peter har en morfar som er polarforsker og er kjent med p
        
        3: Hanne: ønsker å ha mer samarbeid med foreningene, og åpne oss for de forskjellige delene. andre foreninger når andre, så vi skal hjelpe dem inn. styrke internkultur etter 
        
        stabil retning. starte med økonomien
        de nye kommer ikke på kosetirsdag
        
        45: har dere ledet noe før
        Peter - ja men husker det ikke. Det gikk bra for man fø
        
        hanne- leder spf. betaler lønn til utlån. det fungerer.
        
        37: samme som over
        
        41: Hvordan vil dere forholde dere til pampene
        Hanne: Vi skal ikke være en forening som dytter folk ut, men det må være en forståelse for at vi skal ha aktive studenter. Pamper kan være til nytte fordi de kan ting, men nå må vi ha nye. man må gi dem en vennlig dytt ut.
        hvis noen sier det var bedre før kan de dra
        Peter: Det har vært 3. år siden siste ordentlige kompetanseoverføringen. så de har trengtes, men vi må ha nye
        
        22. Hvordan er dere som ledere som personlig
        Peter: Vant til å være i flere foreninger, og vet å sette foreningen først. Etter tillegg fra Inger: gå til sifra, holde det inne
        Hanne: Vil representere hele foreningen, snakke med hovedstyre
        
        47. Vi er veldig glade i å bruke escape. Hvordan vil dere hjelpe med å få fakturaer for utlån sendt først
        Hanne: trengs rutiner
        Peter. Med gamle kassa hentet man ut i slutten av semesteret
        
        6. Hva vil dere si er deres største svakheter
        Peter: ikke så flink til å delegere
        Hanne: Kan være utålmodig. Mendet vet jeg at jeg er, så det kan jeg gjøre noe med.
        
        x5 har dere erfaringer med personalsaker??
        Hanne :ja det har jeg. og det føler jeg at jeg kan
        Peter: har ikke så mye, men learning by doing, og SiO foreninger
        
        9. Internkulturen har vært rar. hvordan vil dere hjelpe å forbedre den
        Peter har ikke så mye på det, for det er ikke bare leder, men HS sammen. trickledown efekt
        
        Hanne: vet ikke hva som menes med rar, men snakke med alle rundt, gå sammen med internansvarlig. og snakke med internmassen.
        
        Takk til kandidatene.
        
        Tellekorpset har et resultat:\\
        \textbf{Peter Hjelle Petersen-Øverleir er valgt som ny leder}
        
        \subsubsection{Nestleder}
        Peter presenterer vervet nestleder.
        Vet ikke helt, har bare gjort masse rart. Hjelper leder.
        To faste ansvar:GF, og protokoll til brreg og kalle in til møter++
        
        Philip Høegh Eriksen stiller:
        En dislaimaer: må løpe så gyda vil svare på spørsmål
        Har endel foreningserfaring, er leder av fadderstyre, sitter i dagen, og startet losjen, og list og buldre-blalbalba
        Har tid fordi jeg bilr ferdig. sosial person, kan bidra med et nettverk blant mellom foreningen og det sosiale som grl og hyggelig person som snakker med mye folk.
        har lyst til å få med en større andel av folket i 3. etasje. og vil være synlig.
        
        ellers energis, blir gira, pusher og drar det litt for langt noen ganger, vet ikke når det skal stoppe, og der
        
        Liker isbjørner mer en pingviner
        
        Har det organisatoriske på plass
        
        16 er du villig til å begynne å like pingviner?
        foretrekker pingviner, men de bor ikke sammen
        er du villig til å ta på deg en pingvindrakt?
        Jeg kan photoshoppes!
        Det blir tatt bilder
        
        46. spørsmålet er trukket.
        
        44. Hva er dine nåværende tanker om kultur og pamper
        Jeg synes pamper er negativt ladet, og alumni er bedre ord. alumni skal tas vare på. og en studentforening skal være for aktive studenter.  kulturen ryddes opp i, ved å dra folk fra andre steder.Erfaringsoverføring er viktig. er du den eneste som kan gjøre noe, så lær det bort.
        
        \textbf{Philip Høegh Eriksen er valgt til nestleder}
        
        \subsubsection{Kasserer}
        Lise orienterer om kasserer vervet.
        Det er ikke så vanskelig som det kanskje virker
        fiks mva, snakk med regnskapsfører. VI har en stabil økogruppa. støtter seg på økonomiansvarlig. Mer økonomiske planleggingen, og melde ting til staten.
        
        Hanne stiller. jeg har litt økonomierfaring fra SPF, har satt meg litt inn i regnskap og slikt
        
        \textbf{Hanne Alexandersdatter Hugem velges til Kasserer.} \textit{(I sak 11.1 er Kasserer vedtatt endret navn til "Økonomiansvarlig")}
    
        \subsubsection{Rekrutteringsansvarlig}
         Adrian orienterer.
        Medlem av hovedstyre
        tar kontakt med de som vil melde seg inn, og spre det gode ord.
        sio foreningsdag, olympiaden, ifi foreningsdag, 
        
        Ingen stiller, vi kommer tilbake til det
        
        \textbf{Etter flere runder så er det ingen som stiller, vervet forblir ledig}


        \subsubsection{Internansvarlig}
        Espen orienterer
        ansvar for intern kultur, og løse saker på en måte s
        
        Vilde Moe stiller. kommer fra uglebo. har erfaring fra internsaker, og er beinhard.
        
        46. Vilde liker du pingviner? ja
        
        45 har du tid js
        
        16 Du nevnte erfaringer fra å håndtere internsaker,  eller skape dem?
        Vilde har erfaring med hvordan internsak prinsipper.
        x2 er du mer effektiv enn meg(tidl. internansvalig)
        
        \textbf{Vilde Moe er valgt til Internansvarlig}
        
    \subsection{Kjellerstyret} %TODO Oppdater listen med verv
        Alle verv som er til valg i kjellerstyret gjelder for ett semester av gangen. Med unntak av Økonomiansvarlig som blir valgt inn for to semester.
        Dag orienterer om KS.
        \subsubsection{Kjellernestleder}
            % \textit{Vervet presenteres.}
        Inger forteller om vervet. Stepper inn for kjellermogul,
        og du gjør med vervet som du vil. 
        
        Fredrik Arman Sarai stiller via Inger. (han vet at han stiller)
        
        11. liker han pingviner
        mest sannsynlig.
        
        \textbf{Fredrik Arman Sarai valgt ved overveldende flertall}
        
        \subsubsection{Barsjef}
            % \textit{Vervet presenteres.}
        Barsjef.
        Elias presenterer barsjef
        hovedansvar for opplæring av sm, fylle opp skift, og leder for bargruppa. ressursperson for folk som står i bar.
        
        Elias stiller. Heter Elias, er nåværende barsjef. vet at jeg klarer det, også fått lært, så nå blir det enda bedre
        
        \textbf{Elias Lynum Ringkjøb valgt ved overveldende flertall}
        
        
        \subsubsection{Kafésjef}
            % \textit{Vervet presenteres.}
        dag presenterer
        fylle skift, vedlikeholde rutiner, følge med på vasking, og være ressursperson for kafegruppa, og gjøre innkjøp til kafe
        
        Inger Marie Liepelt stiller. Vil gå en annen retning i livet. 
        
        14. kommer til du til å finne nye kaffe typer å prøve
        
        45 er du villig til å ha milkshake på kafe menyen?
        
        48. vil du endre traktekaffen
        den som er så jævlig? ja!
        
        \textbf{Inger Marie Liepelt valgt ved overveldene flertall}
        
        \subsubsection{Økonomiansvarlig}
            % \textit{Vervet presenteres.}
        Trym Auren presenterer
        man skal punche tall, føre regnskap
        
        Trym stiller.
        Har hatt vervet i 5 måneder. 
        applaus fra salen.
        1 skal ikke du på utvekslingen?
        ja det skal jeg, så om noen andre vil er
        
        \textbf{Trym Auren valgt til Økonomiansvarlig} \textit{(I sak 11.1 er Økonomiansvarlig vedtatt endret navn til "Kasserer")}

        
        \subsubsection{Innkjøpsansvarlig}
            % \textit{Vervet presenteres.}
        Hilmar orienterer
        
        Mons stiller.
        
        \textbf{Mons Eirik Bjerkan velges til Innkjøpsansvarlig} 
        
        \subsubsection{Teknisk sjef}
            % \textit{Vervet presenteres.}
        kontakt med eiendom, fikse lyspærer, sørge for at vi har teip. Forslag til forbedringer.
        
        settes strek, vi kommer tilbake
        
        Etter flere runder stiller Elias Hoel Birketvedt stiller
        
        14 har du tid
        - ja
        
        \textbf{Elias Hoel Birketvedt valgt til ny Teknisk sjef}
        
        \subsubsection{Utlånsansvarlig}
            % \textit{Vervet presenteres.}
        Hanne presentere
        svare mail, følge opp, sende faktura, ha verv i SPF, sørge for å ikke leie ut når det skjer andre ting
        
        Silje stiller. går 1. år på språktek. er øko i fadderstyret og list, og har verv og
        
        14. kommer du til å sørge for å svare i tide og ikke dobbelbukke
        selvfølgelig.
        
        \textbf{Silje Helgesen velges ved overveldende flertall}
        
        \subsubsection{DJ-sjef}
            % \textit{Vervet presenteres.}
        FInt verv, alle godene, koordinere slik at vi har dj til de store arr. lederansvar for djer. 
        
        \textbf{Ingen stiller til vervet, vervet er ledig}
        
        
        
        \subsubsection{Arrangementskoordinator}

        Jørgen forklarer. er basically arrangementssjef, sørger for at det går greit for det, ressursperson for det praktiske som skal gjennomføres ved arrangement. krever ikke så mye, hvis du har lyst, kan du definere det, eller gjøre hva du vil.
        
        en liten pause
        
        lina stiller. har lyst til å bli bartender. har ikke så mye erfaring, men har lyst til å lære.
        13. hardutidoglikerdupingviner
        Ikke stress og velger gjerne bort ex-phil.
        
        \textbf{Lina Christine Wessel Hauck er valgt til Arrangementskoordinator}

\newpage

\section{Vedtektsforslag}

    \subsection{§3 c, §5 d, §6 d, §6 c og §8 c}
        Endre at navnet på vervet Kasserer i HS og Økonomiansvarlig i KS byttes med hverandre.
        
        Lise presenterer.
        Vi har hatt motsatt av alle andre. nå er vi team normie. hvis noen skal lage drama nå nå får de med meg å gjøre
        
        44 jeg er tidligere økonomiansvarlig, men følte seg som kasserer.
        
        \textbf{Vedtektsendringen er vedtatt ved overveldende flertall}
        
        
    \subsection{§7 a - Skrivefeil}
        Endre:
        \\Generalforsamlingen er foreningens høyeste myndighet. Den er beslutningsdyktig når minst 10 prosent eller 30 stykker av medlemmene er tilstedet.
        \\Til:
        \\Generalforsamlingen er foreningens høyeste myndighet. Den er beslutningsdyktig når minst 10 prosent eller 30 stykker av medlemmene er til \textcolor{ForestGreen}{stede}.
        
        46 kunne vi ikke bare tatt som redakjsonell endring
        martin: ikke opp til meg
        
        \textbf{vedtatt ved overveldende flertall}

    \subsection{§2}
    Legge til nytt punkt (e) i paragraf 2:
    
    \begin{enumerate}
    	\item{Cybernetisk Selskab har tre kategorier medlemmer:}
    	\begin{itemize}
    		\item{Semesterbetalende medlemmer.}
    		\item{Livsvarige medlemmer.}
    		\item{Æresmedlemmer.}
    	\end{itemize}
    	\item{Kontingent for livsvarlig medlemskap er ti ganger ordinær semesterkontingent.}
    	\item{Æresmedlemmer utnevnes på generalforsamling med \sfrac{2}{3} flertall.}
    	\item Medlemmer forplikter seg til å handle i tråd med foreningens formål jf. §1 og \fhref{https://wiki.cyb.no/display/CYB/Retningslinjer}{etiske retningslinjer for forenings- og styremedlemmer} ved instituttet.
    	\item \textcolor{ForestGreen}{Medlemskap kan kun inndras gjennom eksklusjon jf. §10.}
    \end{enumerate}
        vi har en ekslusjonsparagraf. vi har ikke en kontraparagraf som sier at dette er eneste måten man blir ekskludert.
        x2 vi har ikke hatt noen særlig eksklusjobssaker, men det kan bli nødvendig i fremmtiden..
        
        Ingen vil ytre seg i debatten.
        
        \textbf{Vedtektsendringen er vedtatt ved overveldende flertall}

\section{Utdeling av utmerkelser}
Trine legger frem saken.
trine forklarer poengene.
skal dele ut 3 pins.
x-sjef kamil skal få en pin.
@hilmar
kjellerpotet inger

lise får gullkort.


\section{Eventuelt}
Gyda får blomster og takketale.

Møtet heves


\newpage




\textbf{\huge{Hovedstyret består nå av:}}

\mbox{}


\begin{center}
\setstretch{2.0}
    \LARGE Leder -- Peter Hjelle Petersen-Øverleir
    \\ Nestleder -- Philip Høegh Eriksen
    \\ Økonomiansvarlig -- Hanne Alexandersdatter Hugem
    \\ Kjellermogul -- Dag Aleksander Hagebø Sommer
    \\ Arrangementssjef -- Robin Simonsen
    \\ Rekrutteringsansvarlig -- Ledig
    \\ Internansvarlig -- Vilde Moe
    \\ Promoteringssjef -- Ida Eneroth Klundseter
    \\ X-sjef -- Kamil K\MakeLowercase{RZYŻANOWSKI}
\end{center}

\mbox{}
\\[24pt]

\begin{center}
\LARGE Protokollunderskrivere:\\[30pt]

\begin{multicols}{2}
    \Large 
    Martin Kristoffer Jørstad Heggem\\[48pt]
    \rule{220pt}{1pt}
    
    Kenneth André Sydness\\[48pt]
    \rule{220pt}{1pt}
    
    
\end{multicols}
\end{center}



\end{document}
